\documentclass[11pt,a4paper,parskip=half-]{scrartcl}

%-------------------------
% Packages
%-------------------------
\usepackage[utf8]{inputenc}
\usepackage{lmodern}
\usepackage{microtype}

\usepackage{amsmath, amssymb, amsfonts, mathtools}
\usepackage{amsthm}

\usepackage{algorithm}
\usepackage{algpseudocode}
\usepackage{minted}

\usepackage{fnpct}
\usepackage[onehalfspacing]{setspace}
\usepackage{enumitem}

\usepackage{geometry}
\geometry{margin=1in}

\usepackage{authoraftertitle}

\usepackage[breaklinks]{hyperref}
\usepackage{cleveref}

% make footnotes rules as wide as the text
\renewcommand{\footnoterule}{%
  \kern -4pt
  \hrule width \textwidth height 0.4pt
  \kern 3.6pt
}

% define helper command for typesetting inline code
\newcommand{\code}[1]{\texttt{#1}}

\setstretch{1.125}

%-------------------------
% Theorem-like environments
%-------------------------
\theoremstyle{definition}
\newtheorem{definition}{Definition}[section]
\newtheorem{example}{Example}[section]

\theoremstyle{plain}
\newtheorem{theorem}{Theorem}[section]
\newtheorem{lemma}{Lemma}[section]
\newtheorem{proposition}{Proposition}[section]

\theoremstyle{remark}
\newtheorem{remark}{Remark}[section]

%-------------------------
% Common notation macros
%-------------------------
\newcommand{\R}{\mathbb{R}}
\newcommand{\N}{\mathbb{N}}
\newcommand{\norm}[1]{\left\lVert #1\right\rVert}
\newcommand{\ip}[2]{\left\langle #1,#2\right\rangle}

% Tensor Train notation
\newcommand{\TT}{\mathrm{TT}}
\newcommand{\rank}{\text{rank}}
\newcommand{\diag}{\mathrm{diag}}

% Optional: consistent algorithm keyword styling
\algrenewcommand\algorithmicrequire{\textbf{Input:}}
\algrenewcommand\algorithmicensure{\textbf{Output:}}

%-------------------------
% Title info
%-------------------------
\title{Tensor Train (TT) Decomposition and PDE Solvers}
\author{Pavel Altmann, Jakub Kopčil, Eva Trejbalová (WREN)}

\hypersetup{
  pdfborderstyle={/S/U/W 0.5},
  pdftitle={\MyTitle},
  pdfauthor={\MyAuthor}
}
\urlstyle{same}

\begin{document}
\maketitle

\begin{abstract}
  This document collects definitions and core algorithms for the tensor train
  (TT) format, basic operations in TT, and how TT representations can be used to
  solve higher-dimensional (4--10 dimensions) partial differential equations
  (PDEs) in Python using NumPy.
\end{abstract}

\newpage
\tableofcontents
\newpage

%==================================================
\section{Motivation}
%==================================================

The goal of the Tensor Train (TT) decomposition is to represent
such tensors using a number of parameters that grows only
linearly in $d$, provided certain low-rank structure exists.

High–dimensional problems naturally lead to tensors, i.e.,
multidimensional arrays of numbers.

If a function
\[
  f(x_1,\dots,x_d)
\]
is sampled on a tensor grid with $n_k$ points in each dimension,
the resulting data structure is a tensor
\[
  A \in \mathbb{R}^{n_1 \times \cdots \times n_d}.
\]

Storing such a full tensor requires
\[
  n_1 n_2 \cdots n_d
\]
numbers, which grows exponentially in $d$. This phenomenon is known as the
\textbf{curse of dimensionality} \cite{Oseledets2009TreeTucker,Oseledets2011TT}.
Low–rank tensor formats aim to overcome this exponential growth.

%==================================================
\section{Tensor Definition}
%==================================================

\begin{definition}[Tensor]
  A $d$-dimensional tensor $A$ is a multidimensional array
  \[
    A \in \mathbb{R}^{n_1 \times n_2 \times \cdots \times n_d},
  \]
  whose entries are indexed as
  \[
    A(i_1,\dots,i_d), \quad 1 \le i_k \le n_k.
  \]
\end{definition}

\paragraph{Explanation (informal).}

A tensor is simply a generalization of vectors and matrices to dimensions $d >
2$. For example, a video can be thought of as a tensor with 4 dimensions: width
$\times$ height $\times$ color channels $\times$ time. The problem is that
storing all entries quickly becomes impossible when $d$ grows.

\paragraph{Representation in code.}

In numerical computing, a tensor is typically stored as an \code{ndarray}
(NumPy array).

\begin{definition}[Mode-$k$ unfolding]
  Let $A \in \mathbb{R}^{n_1 \times \cdots \times n_d}$.

  The $k$-th unfolding (or matricization) of $A$ is the matrix
  \[
    A^{\langle k\rangle}
    \in
    \mathbb{R}^{(n_1 \cdots n_k)\times(n_{k+1}\cdots n_d)}
  \]
  obtained by reshaping the tensor so that the first $k$
  indices form the row index and the remaining indices form
  the column index.
\end{definition}

\paragraph{Explanation (informal).}

Unfolding allows us to covert a tensor with arbitrary dimensions into a single
matrix by reshaping (similarly to how a matrix can be represented by a regular
array given proper indexing). This is important because matrices can be
approximated by low-rank factorizations using the Singular Value Decomposition
(SVD).

The SVD decomposition is of the form
\[
  A = U \Sigma V^\top,
\]
and truncating small singular values yields the best low-rank approximation in
the Frobenius norm.

TT decomposition is built from successive low-rank approximations of unfolding
matrices\cite{Oseledets2011TT}.

\paragraph{Representation in code.}

In NumPy, unfolding is just a reshape:

\begin{minted}{python}
A_k = A.reshape(
    np.prod(A.shape[:k]),
    np.prod(A.shape[k:])
)
\end{minted}

%==================================================
\section{Tensor Train Decomposition}
%==================================================

\subsection{Definition}
\begin{definition}[Tensor Train decomposition (TTD)]
  A tensor $A \in \mathbb{R}^{n_1 \times \cdots \times n_d}$
  is in Tensor Train (TT) format if it can be written as
  \[
    A(i_1,\dots,i_d)
    =
    G_1(i_1) G_2(i_2) \cdots G_d(i_d),
  \]
  where for each $k$:
  \[
    G_k(i_k) \in \mathbb{R}^{r_{k-1} \times r_k},
  \]
  and
  \[
    r_0 = r_d = 1.
  \]

  Equivalently, treating matrices as a function \(\R^2\to \R\):
  \[
    A(i_1,\dots,i_d)
    =
    \sum_{\alpha_1,\dots,\alpha_{d-1}}
    G_1(i_1)(1,\alpha_1)
    G_2(i_2)(\alpha_1,\alpha_2)
    \cdots
    G_d(i_d)(\alpha_{d-1},1).
  \]
\end{definition}

\begin{remark}
  The storage complexity of the TT format is
  \[
    \sum_{k=1}^d n_k r_{k-1} r_k.
  \]

  If $r_k \le r$ and $n_k = n$,
  this becomes
  \[
    \mathcal{O}(d n r^2),
  \]
  which is linear in $d$. Thus, TT format breaks the curse of dimensionality
  when TT-ranks remain moderate.
\end{remark}

\paragraph{Explanation (informal).}
Each index $i_k$ selects a matrix $G_k(i_k)$. Multiplying these matrices yields
a scalar because the boundary ranks satisfy $r_0 = r_d = 1$.

Thus, the tensor entry is obtained by a chain of matrix multiplications.

\paragraph{Explanation (informal).}
Instead of storing the full tensor, we store small 3D tensors called
\textbf{cores}. Each core connects only to its neighbors, forming a chain (a
``train''). Instead of storing all $n^d$ elements, only approximately
$\mathcal{O}(d n r^2)$ numbers are required when ranks $r_k$ are
moderate\cite{Oseledets2011TT}.

\paragraph{Example (TT representation).}
Let $A \in \mathbb{R}^{2 \times 2 \times 2}$ be defined by
\[
A(i_1,i_2,i_3) = i_1 + i_2 + i_3.
\]

A TT representation writes
\[
A(i_1,i_2,i_3) = G_1(i_1)\,G_2(i_2)\,G_3(i_3),
\]
with
\[
G_1(i_1) \in \mathbb{R}^{1 \times r_1}, \quad
G_2(i_2) \in \mathbb{R}^{r_1 \times r_2}, \quad
G_3(i_3) \in \mathbb{R}^{r_2 \times 1}.
\]

\paragraph{Explanation (informal).}
The tensor is not stored entrywise. Instead, each index selects a small matrix,
and multiplying these matrices produces the scalar value.

\paragraph{Explanation (informal).}
For structured tensors (such as sums of variables), small ranks $r_k$ are
sufficient. This yields a compressed representation.

\paragraph{Representation in code.}

In practice, TT format is stored as a list of cores:

\begin{minted}{python}
# TT tensor represented as list of cores
G = [G1, G2, ..., Gd]
# Each core has shape (r_{k-1}, n_k, r_k)
Gk.shape = (r_prev, n_k, r_next)
\end{minted}

\subsection{TT-rank and related concepts}

\begin{definition}[TT-rank]
  Let $A^{\langle k\rangle}$ denote the $k$-th unfolding of $A$.
  The TT-ranks are defined as
  %TODO: better lay explenation
  \[
    r_k = \operatorname{rank}\!\left(A^{\langle k\rangle}\right),
    \quad k=1,\dots,d-1.
  \]
\end{definition}

\paragraph{Explanation (informal).}
The rank $r_k$ measures how strongly the first $k$ variables interact with the
remaining $d - k$ variables. If $r_k = 1$, the first $k$ variables are
independent of the rest.

These ranks determine the minimal possible internal dimensions of any TT
representation.

%==================================================
\section{Basic Operations in TT Format}
%==================================================

\subsection{Computing TTD}
\subsubsection{From Unfoldings to TT-SVD}
\begin{definition}[Frobenius norm of tensor]
  For a tensor $A$,
  \[
    \|A\|_F
    =
    \sqrt{
      \sum_{i_1,\dots,i_d}
      A(i_1,\dots,i_d)^2
    }.
  \]
\end{definition}

\paragraph{Conceptual idea.}
The TT-SVD algorithm constructs a TT representation
by applying truncated SVD to successive unfolding matrices.

At each step:

\begin{enumerate}
  \item Reshape the tensor into a matrix.
  \item Compute its truncated SVD.
  \item Store left singular vectors as a TT core.
  \item Pass the remainder to the next step.
\end{enumerate}

This produces a quasi-optimal low-rank TT approximation.

\subsubsection{Algorithm for computing a TT decomposition (TT-SVD)}

The TT-SVD algorithm constructs matrix-valued functions
\[
  G_k : \{1,\dots,n_k\} \to \mathbb{R}^{r_{k-1}\times r_k}
\]
by sequential low-rank approximation of unfolding matrices
\cite{Oseledets2011TT} as described in \cref{alg:ttsvd}.

\begin{algorithm}[ht]
  \caption{TT-SVD in matrix-valued function notation}
  \label{alg:ttsvd}
  \begin{algorithmic}[1]
    \Require Tensor $A \in \mathbb{R}^{n_1 \times \cdots \times n_d}$, tolerance
    $\varepsilon$
    \Ensure Matrix-valued functions $G_1,\dots,G_d$

    \State Set $r_0 \gets 1$
    \State Set $C \gets A$
    \State Set $\delta \gets \dfrac{\varepsilon}{\sqrt{d-1}} \|A\|_F$

    \For{$k=1$ to $d-1$}

    \State Reshape $C$ into matrix
    \[
      C \in
      \mathbb{R}^{(r_{k-1} n_k) \times (n_{k+1}\cdots n_d)}
    \]

    \State Compute truncated SVD
    \[
      C \approx U \Sigma V^\top,
      \quad
      \|C-U\Sigma V^\top\|_F \le \delta
    \]

    \State Let $r_k$ be the number of kept singular values

    \State Partition $U$ into $n_k$ consecutive blocks of size
    $r_{k-1} \times r_k$

    \State Define matrix-valued function
    \[
      G_k(i_k) :=
      \text{the block of $U$ corresponding to index $i_k$}
    \]

    \State Set $C \gets \Sigma V^\top$

    \EndFor

    \State Define $G_d(i_d)$ from the final matrix $C$
    \[
      C \in \mathbb{R}^{r_{d-1} \times n_d}
    \]
    by taking column $i_d$ as a matrix in $\mathbb{R}^{r_{d-1}\times 1}$.

  \end{algorithmic}
\end{algorithm}

\begin{remark}
  A standard relative-accuracy choice is to distribute the Frobenius error
  budget across the $d-1$ SVD truncations:
  \[
    \delta = \frac{\varepsilon}{\sqrt{d-1}}\,\|A\|_F,
  \]
  so that the final TT approximation satisfies $\|A-\widehat{A}\|_F
  \le \varepsilon\|A\|_F$
  (see the TT-SVD error bound in \cite{Oseledets2011TT}).
\end{remark}

\paragraph{Explanation (informal).}
At step $k$ we reshape the tensor into a matrix that separates
the first $k$ indices from the remaining $d-k$ indices.
This matrix describes how these two groups of indices
depend on each other.

We then compute a truncated Singular Value Decomposition (SVD)
of this matrix to obtain a low-rank approximation.
The rank $r_k$ of this approximation determines the next
TT-rank.

The left singular vectors are reshaped into blocks that form
the $k$-th TT core $G_k(i_k)$.
The remaining part of the tensor is passed to the next step,
where the same procedure is applied again.

By repeating this process for $k = 1,\dots,d-1$,
the algorithm constructs all TT cores and determines
the TT-ranks of the representation.

In this way, TT-SVD builds the tensor train representation
sequentially, extracting one core at a time.

\paragraph{Example (first TT-SVD step).}
Let $A \in \mathbb{R}^{2 \times 3 \times 2}$.

At step $k = 1$, reshape:
\[
A^{\langle 1 \rangle} \in \mathbb{R}^{2 \times 6}.
\]

Compute truncated SVD:
\[
A^{\langle 1 \rangle} \approx U \Sigma V^\top.
\]

Let $r_1$ be the number of retained singular values.

Define:
\[
G_1(i_1) = \text{block of } U \text{ corresponding to } i_1,
\]
and set
\[
C = \Sigma V^\top.
\]

\paragraph{Explanation (informal).}
We separate the first variable from the remaining ones and approximate their
interaction by a low-rank matrix.

\paragraph{Explanation (informal).}
The matrix $U$ encodes how $i_1$ interacts with the rest. The remainder $C$
is passed to the next step.

\paragraph{Code slot (NumPy).}
\begin{minted}{python}
  # TODO implemented TTD
\end{minted}

\subsection{Rounding}

\subsubsection{Algorithm for rounding (rank truncation / recompression)}

Rounding (also called recompression) reduces inflated TT-ranks after operations
(addition, etc.) while controlling the approximation error. The standard
procedure is to orthogonalize (QR sweep) and then truncate (SVD sweep)
\cite{Oseledets2011TT} as shown in \cref{alg:rounding}.

\begin{algorithm}[ht]
  \caption{TT rounding in matrix-valued function notation}
  \label{alg:rounding}
  \begin{algorithmic}[1]
    \Require Matrix-valued functions $G_1,\dots,G_d$, tolerance $\varepsilon$
    \Ensure Rounded matrix-valued functions
    $\widetilde{G}_1,\dots,\widetilde{G}_d$

    \State Set $\delta \gets \dfrac{\varepsilon}{\sqrt{d-1}} \|\widehat{A}\|_F$

    \State \textbf{Left-to-right orthogonalization}

    \For{$k=1$ to $d-1$}

    \State Form matrix
    \[
      M_k :=
      \begin{bmatrix}
        G_k(1) \\
        \vdots \\
        G_k(n_k)
      \end{bmatrix}
      \in
      \mathbb{R}^{(r_{k-1} n_k) \times r_k}
    \]

    \State Compute thin QR:
    \[
      M_k = Q_k R_k
    \]

    \State Redefine $G_k(i_k)$ as corresponding block of $Q_k$

    \State Update next core:
    \[
      G_{k+1}(i_{k+1})
      \leftarrow
      R_k \, G_{k+1}(i_{k+1})
    \]

    \EndFor

    \State \textbf{Right-to-left truncation}

    \For{$k=d$ down to $2$}

    \State Form matrix
    \[
      N_k :=
      \begin{bmatrix}
        G_k(1) & \cdots & G_k(n_k)
      \end{bmatrix}
      \in
      \mathbb{R}^{r_{k-1} \times (n_k r_k)}
    \]

    \State Compute truncated SVD
    \[
      N_k \approx U \Sigma V^\top
    \]

    \State Update:
    \[
      G_k(i_k)
      \leftarrow
      \text{corresponding block of } V^\top
    \]

    \State Update previous core:
    \[
      G_{k-1}(i_{k-1})
      \leftarrow
      G_{k-1}(i_{k-1}) (U\Sigma)
    \]

    \EndFor

  \end{algorithmic}
\end{algorithm}

\begin{remark}
  The QR sweep makes the partial products
  \[
    G_1(i_1)\cdots G_k(i_k)
  \]
  orthonormal in Frobenius norm.

  This ensures that the subsequent SVD truncations behave like
  optimal matrix approximations, leading to numerical stability
  and quasi-optimal error control.
\end{remark}

\paragraph{Important distinction.}
\begin{itemize}
  \item TT-SVD constructs a TT representation from a full tensor.
  \item Rounding reduces the ranks of an existing TT tensor
    while controlling the approximation error.
\end{itemize}

\paragraph{Example (rank growth and truncation).}
Let $A$ and $B$ be TT tensors with ranks
\[
r^A = (1,3,3,1), \quad r^B = (1,4,4,1).
\]

After addition:
\[
r^C = (1,7,7,1).
\]

After rounding:
\[
r^C \rightarrow (1,3,3,1).
\]

\paragraph{Explanation (informal).}
Basic operations increase TT-ranks. This leads to higher storage and
computational cost.

\paragraph{Explanation (informal).}
Rounding reduces ranks while controlling the approximation error. It restores
a compact representation.

\paragraph{Note.}
The Frobenius norm satisfies $\|A\|_F = \sqrt{\langle A,A\rangle}$, so it
reduces to an inner product computation (see \cref{sec:inner_product}).

\paragraph{Code slot (NumPy).}
\begin{minted}{python}
# TODO implemented rounding
\end{minted}

\subsection{Addition}

Given $A$ and $B$ in TT format, a standard construction forms a TT
representation of
$C=A+B$ by stacking cores blockwise; ranks add (then we round).

\paragraph{Construction (conceptual).}
If $A$ has cores $G_A^{(k)} \in \mathbb{R}^{r_{k-1}^A \times n_k \times r_k^A}$
and $B$ has $G_B^{(k)} \in \mathbb{R}^{r_{k-1}^B \times n_k \times r_k^B}$,
define cores of $C$:

\begin{itemize}[leftmargin=2em]
  \item For $k=1$: concatenate along the last (right-rank) index:
    \[
      G_C^{(1)}(1,i_1,:) = \big[\,G_A^{(1)}(1,i_1,:)
      \ \ \ G_B^{(1)}(1,i_1,:)\,\big]
    \]
  \item For $1<k<d$: block-diagonal stacking:
    \[
      G_C^{(k)}(:,i_k,:) =
      \begin{bmatrix}
        G_A^{(k)}(:,i_k,:) & 0\\
        0 & G_B^{(k)}(:,i_k,:)
      \end{bmatrix}
    \]
  \item For $k=d$: concatenate along the first (left-rank) index:
    \[
      G_C^{(d)}(:,i_d,1) =
      \begin{bmatrix}
        G_A^{(d)}(:,i_d,1)\\
        G_B^{(d)}(:,i_d,1)
      \end{bmatrix}
    \]
\end{itemize}

This yields ranks $r_k^C = r_k^A + r_k^B$. The procedure is finalized by
rounding $C$ (\cref{alg:rounding}).

\paragraph{Code slot.}
\begin{minted}{python}
# then call tt_round(...)
\end{minted}

\paragraph{Scalar multiplication.}
For $c\in\mathbb{R}$ and TT tensor $A$ with cores $G^{(k)}$, one may
scale a single core:
\[
  (cA)\ \text{can be represented by cores}\ (c\,G^{(1)}),G^{(2)},\dots,G^{(d)}.
\]

\paragraph{Code slot.}
\begin{minted}{python}
# - scalar multiply: scale G[0] *= c
\end{minted}

\paragraph{Code slot.}
\begin{minted}{python}
  # TODO implemented proper norm_F
\end{minted}

% TODO: Norm computation via contractions; mention orthogonalization.

\subsection{Inner products}
\label{sec:inner_product}

\begin{definition}[Inner product in TT format]

  Let $A$ and $B$ be two tensors in TT format with
  matrix-valued cores

  \[
    G_k : \{1,\dots,n_k\} \to \mathbb{R}^{r_{k-1}\times r_k},
    \quad
    H_k : \{1,\dots,n_k\} \to \mathbb{R}^{s_{k-1}\times s_k}.
  \]

  The Frobenius inner product

  \[
    \langle A,B \rangle
    =
    \sum_{i_1,\dots,i_d}
    A(i_1,\dots,i_d)\,
    B(i_1,\dots,i_d)
  \]

  can be computed sequentially by summing the inner products of
  the cores contracted by matrices $M_k$ as shown in \cref{alg:inner_product}.

  Define $M_0 = [1]$ and recursively compute

  \[
    M_k
    =
    \sum_{i_k=1}^{n_k}
    G_k(i_k)^\top
    \, M_{k-1} \,
    H_k(i_k),
    \qquad k=1,\dots,d.
  \]

  Then

  \[
    \langle A,B \rangle = M_d.
  \]

  \begin{algorithm}[t]
    \caption{Inner product in TT format}
    \label{alg:inner_product}
    \begin{algorithmic}[1]
      \Require Matrix-valued functions $G_1,\dots,G_d$

      \Ensure Inner product $\langle A,B \rangle$

      \State Set $M_0 \coloneqq [1]$ \Comment{$1 \times 1$ matrix}

      \For{$k = 1$ \textbf{to} $d$}

      \State $M_k = \sum_{i_k} G_k(i_k)^T M_{k-1} H_k(i_k)$
      \Comment{$M_k$ is $r_k \times r_k$}

      \EndFor

      \State \Return $M_d$ \Comment{$M_d$ is a $1 \times 1$ matrix}
    \end{algorithmic}
  \end{algorithm}
\end{definition}

\paragraph{Explanation (informal).}
Instead of summing over all $n_1\cdots n_d$ tensor entries,
we contract the tensor cores one dimension at a time.

At each step $k$, we accumulate a small matrix $M_k$ that carries the
interaction between the first $k$ dimensions. Since the boundary ranks satisfy
$r_0=s_0=1$, the final matrix $M_d$ has size $1\times 1$ and thus can be
interpreted as a scalar. This reduces exponential complexity to polynomial
complexity.

\paragraph{Complexity.}

If $n_k \le n$ and all TT-ranks are bounded by $r$, the computational complexity
reduces from exponential $\mathcal{O}(n^d)$ to polynomial $\mathcal{O}(dnr^3)$.

\paragraph{Code slot.}
\begin{minted}{python}
  # inner product
\end{minted}

\subsection{Flipping an axis}

\paragraph{Definition (axis flip).}
Let $A \in \mathbb{R}^{n_1 \times \cdots \times n_d}$ be in TT format with cores
\[
G_k \in \mathbb{R}^{r_{k-1} \times n_k \times r_k}.
\]

Flipping axis $k$ produces a tensor $\tilde{A}$ such that
\[
\tilde{A}(i_1,\dots,i_k,\dots,i_d)
=
A(i_1,\dots,n_k - i_k + 1,\dots,i_d).
\]

\paragraph{Explanation (informal).}
Flipping an axis corresponds to reversing the ordering of indices along that
dimension.

Since each dimension is represented independently in TT format, the operation
affects only a single core.

\paragraph{Construction.}
The flipped tensor is obtained by reversing the index dimension of the $k$-th core:
\[
\tilde{G}_k(:, i_k, :) = G_k(:,\, n_k - i_k + 1, :).
\]

All other cores remain unchanged:
\[
\tilde{G}_j = G_j, \quad j \neq k.
\]

\paragraph{Remark.}
The TT-ranks are unchanged by this operation.

\paragraph{Complexity.}
The operation requires $\mathcal{O}(r_{k-1} n_k r_k)$ operations.

\paragraph{Code slot (NumPy).}
\begin{minted}{python}
#TODO implement
\end{minted}

%==================================================
\subsection{Element access}
%==================================================

\begin{definition}[Element evaluation]

Let $A$ be a TT tensor with matrix-valued cores
\[
G_k : \{1,\dots,n_k\} \to \mathbb{R}^{r_{k-1}\times r_k}.
\]

The entry
\[
A(i_1,\dots,i_d)
\]
can be computed by sequential multiplication:
\[
v_0 = 1,
\quad
v_k = v_{k-1} G_k(i_k),
\quad
k=1,\dots,d.
\]

Then
\[
A(i_1,\dots,i_d) = v_d.
\]

\end{definition}

\paragraph{Explanation (informal).}

Instead of reconstructing the full tensor, we select one slice from each core
and multiply them from left to right.

Each step propagates a small matrix (or vector), and since $r_0=r_d=1$,
the final result is a scalar.

\paragraph{Complexity.}

If TT-ranks are bounded by $r$, the cost is
\[
\mathcal{O}(d r^2).
\]

\paragraph{Code slot.}
\begin{minted}{python}
# TODO: element access
\end{minted}

%==================================================
\subsection{Mode permutation (transpose)}
\label{sec:mode-permutation}
%==================================================

\begin{definition}[Mode permutation in TT format]

Let $A \in \mathbb{R}^{n_1 \times \cdots \times n_d}$ be represented
in TT format with cores $G_1,\dots,G_d$.

Given a permutation
\[
\sigma : \{1,\dots,d\} \to \{1,\dots,d\},
\]
the permuted tensor $\widetilde{A}$ is defined by
\[
\widetilde{A}(i_1,\dots,i_d)
=
A(i_{\sigma(1)},\dots,i_{\sigma(d)}).
\]

A TT representation of $\widetilde{A}$ is obtained by
reordering and recompressing the TT cores such that the
tensor remains in TT format.

\end{definition}

\paragraph{Explanation (informal).}

Unlike basic operations such as scaling or axis flipping,
a general permutation cannot be performed by simply reordering
the TT cores.

This is because each core encodes dependencies between
neighboring dimensions.

Changing the order of dimensions changes these dependencies,
so the TT decomposition must be recomputed locally.

\paragraph{Construction (adjacent swap).}

A general permutation is implemented as a sequence of swaps
of neighboring dimensions.

To swap dimensions $k$ and $k+1$:

\begin{enumerate}
  \item Contract the two neighboring cores:
  \[
  T(i_k,i_{k+1})
  =
  G_k(i_k) \, G_{k+1}(i_{k+1})
  \]

  \item Reshape $T$ into a matrix:
  \[
  T \in \mathbb{R}^{(r_{k-1} n_{k+1}) \times (n_k r_{k+1})}
  \]

  \item Compute a truncated SVD:
  \[
  T \approx U \Sigma V^\top
  \]

  \item Reshape $U$ and $V^\top$ into new TT cores:
  \[
  \widetilde{G}_k(i_{k+1}), \quad \widetilde{G}_{k+1}(i_k)
  \]

\end{enumerate}

\paragraph{Remark.}

Mode permutation generally increases TT-ranks before truncation.
Therefore, rounding (\cref{alg:rounding}) is required to maintain
a compact representation.

\paragraph{Complexity.}

A single adjacent swap has complexity
\[
\mathcal{O}(n_k n_{k+1} r^3).
\]

A general permutation requires $\mathcal{O}(d^2)$ such swaps.

\paragraph{Code slot.}
\begin{minted}{python}
# TODO: implement adjacent swap of TT cores
# then compose swaps to achieve general permutation
\end{minted}

%==================================================
\subsection{Tensor contraction (tensordot)}
%==================================================

\begin{definition}[Tensor contraction in TT format]

Let $A$ and $B$ be TT tensors.

A tensor contraction (tensordot) along selected modes sums over
shared indices of $A$ and $B$.

In the special case where all modes are contracted, this reduces to
the Frobenius inner product:
\[
\langle A,B\rangle =
\sum_{i_1,\dots,i_d}
A(i_1,\dots,i_d)\,B(i_1,\dots,i_d).
\]

\end{definition}

\paragraph{Remark (ordering of contracted modes).}

Efficient contraction in TT format requires the contracted modes
to form a contiguous block of TT cores.

More precisely, the contracted cores should appear consecutively
and include a boundary core:
\begin{itemize}
  \item either starting from the first core,
  \item or ending at the last core.
\end{itemize}

If the required contraction does not satisfy this ordering,
the TT cores must first be permuted (see \cref{sec:mode-permutation})
before performing the contraction.

\paragraph{Explanation (informal).}

The contraction is performed sequentially across dimensions.

At each dimension $k$, we combine slices of $A$ and $B$,
sum over the contracted index, and propagate a small matrix
that carries the accumulated interaction between all previously
contracted dimensions.

More precisely, if $M_{k-1}$ denotes the propagated matrix from
the previous step, the next contraction has the form
\[
M_k
=
\sum_{i_k}
G_k(i_k)^\top
\, M_{k-1} \,
H_k(i_k),
\]
which reduces the contribution of dimension $k$ into a new
small matrix $M_k$.

This avoids forming the full tensors explicitly.

\paragraph{Important note.}

General contractions increase TT-ranks. This is mostly due to the fact that we need to permutate the cores before performing the tensordot.
Therefore, rounding (\cref{alg:rounding}) is applied afterward.

\paragraph{Complexity.}

If $n_k \le n$ and ranks are bounded by $r$, the cost is
\[
\mathcal{O}(d n r^3).
\]

\paragraph{Code slot.}
\begin{minted}{python}
# TODO: tensordot / contraction
\end{minted}

%==================================================
\subsection{Differential Operators}
\label{sec:differential_operators}
%==================================================

Classical tensor-train approaches often represent differential
operators as TT-matrices and apply them through TT-matrix--TT-vector
multiplication.

In this work, however, differential operators are applied directly
to TT cores. This avoids constructing explicit TT-matrix
representations and operates directly on the TTD cores.

%==================================================
\subsubsection{Gradient in TTD}
\label{sec:gradient_ttd}
%==================================================

\begin{definition}[Gradient in TTD format]

Let
\[
A(i_1,\dots,i_d)
=
G_1(i_1)\cdots G_d(i_d)
\]
be a tensor represented in TT format.

In our implementation, the partial derivative along axis $k$
is approximated by applying a finite-difference rule directly
to the physical index of the $k$-th TT core while leaving all
other cores unchanged.

The resulting approximation is

\[
\frac{\partial A}{\partial x_k}(i_1,\dots,i_d)
\approx
G_1(i_1)\cdots \widetilde{G}_k(i_k)\cdots G_d(i_d),
\]

where

\[
\widetilde{G}_k
=
\operatorname{FD}(G_k),
\]

and $\operatorname{FD}$ denotes a finite-difference approximation
along the physical index of the core.

\end{definition}

\paragraph{Explanation (informal).}

The TT representation stores each spatial dimension in the physical
index of a separate core. This locality allows directional derivatives
to be computed without modifying the entire tensor representation.

Only the core associated with the selected spatial dimension must be
updated, while all remaining cores are reused.

\paragraph{Implementation.}

For a selected axis $k$:

\begin{enumerate}
  \item Copy all TT cores.
  \item Apply a finite-difference stencil to the physical dimension
        of core $G_k$.
  \item Replace $G_k$ by the differentiated core
        $\widetilde{G}_k$.
  \item Construct a new TTD from the modified cores.
\end{enumerate}

If derivatives are requested along multiple axes,
the procedure is repeated independently for each axis.

\paragraph{Numerical note.}

If the variation of a core along its physical dimension
falls below a prescribed tolerance, the derivative core
is replaced by zero to avoid amplification of numerical noise.

\paragraph{Output convention.}

If a single axis is specified, the result is a single TTD.

If multiple axes are specified, the result is a tuple
containing one TTD for each derivative direction.

%==================================================
\subsubsection{Divergence}
\label{sec:divergence_ttd}
%==================================================

\begin{definition}[Divergence]

Let
\[
F = (F_1,\dots,F_d)
\]
be a vector field whose components are represented in TTD format.

The divergence is defined as
\[
\nabla \cdot F
=
\sum_{k=1}^d
\frac{\partial F_k}{\partial x_k}.
\]

Each partial derivative is computed using the TTD gradient operation,
and the resulting TTD tensors are added together.

\end{definition}

\paragraph{Explanation (informal).}

The divergence is computed by differentiating each vector-field component
along its matching coordinate direction and summing the results.

No explicit differential operator is assembled.

\paragraph{Implementation choice.}

In our implementation, the divergence is obtained by computing the
gradient of the vector field and extracting the trace of the resulting
Jacobian tensor.

\paragraph{Code slot.}
\begin{minted}{python}
# TODO: divergence in TTD
# - compute gradient of component F_k along axis k
# - sum all derivative TTDs
# - round if needed
\end{minted} 

%==================================================
\subsubsection{Laplace operator}
%==================================================

\begin{definition}[Laplace operator]

For a scalar tensor $A$, the Laplace operator is computed as
\[
\Delta A
=
\nabla \cdot (\nabla A)
=
\sum_{k=1}^d
\frac{\partial^2 A}{\partial x_k^2}.
\]

In TTD format, this is obtained by applying the gradient operation twice:
first to compute directional derivatives, then again to differentiate each
directional derivative along the same axis.

\end{definition}

\paragraph{Explanation (informal).}

The Laplace operator is implemented as divergence of the gradient.

\paragraph{Important note.}

The summation of the second-derivative TTDs may increase TT-ranks.
Therefore, rounding should be applied after the sum.

\paragraph{Code slot.}
\begin{minted}{python}
# TODO: Laplace in TTD
# - compute gradient(A)
# - compute divergence(gradient(A))
# - round result
\end{minted}

%==================================================
\section{PDEs and Using TT format to solve them}
%==================================================

\subsection{Definition of PDE}

\begin{definition}[Partial Differential Equation]

  Let $\Omega \subset \mathbb{R}^d$.
  A partial differential equation (PDE) has the form

  \[
    \mathcal{L} u(x_1,\dots,x_d) = f(x_1,\dots,x_d),
  \]

  where

  \begin{itemize}
    \item $u : \Omega \to \mathbb{R}$ is the unknown function,
    \item $\mathcal{L}$ is a differential operator,
    \item $f$ is a given function.
  \end{itemize}

  Example (Poisson equation):

  \[
    -\Delta u = f,
    \qquad
    \Delta = \sum_{k=1}^d \frac{\partial^2}{\partial x_k^2}.
  \]

\end{definition}

\paragraph{Explanation (informal).}

A PDE describes how a function of several variables changes in space.

For example, the Poisson equation models heat distribution, electrostatics, or
diffusion.

In low dimensions ($d=1,2,3$), standard numerical methods work well.
However, when $d$ becomes large (for example $d=6,8,10$), the computational cost
grows rapidly.

Tensor Train methods allow us to represent the solution $u$ in compressed
TT form, while differential operations are approximated directly on the TT
cores.

\subsection{Discretization and the Curse of Dimensionality}

\begin{definition}[Finite-dimensional discretization]

  Discretize each variable $x_k$ using $n$ grid points.

  The continuous function $u(x_1,\dots,x_d)$ is replaced by a tensor

  \[
    U \in \mathbb{R}^{n \times \cdots \times n}
    = \mathbb{R}^{n^d}.
  \]

  The differential operator $\mathcal{L}$ is replaced by a matrix

  \[
    A \in \mathbb{R}^{n^d \times n^d}.
  \]

  The PDE becomes the linear system

  \[
    A U = F.
  \]

\end{definition}

\paragraph{Explanation (informal).}

After discretization, solving a PDE means solving a linear system. If each
dimension has $n$ grid points, the number of unknowns is $n^d$.

Even for moderate values, for example $n=50$ and $d=8$, the system has $50^8
\approx 4 \times 10^{13}$ unknowns — far beyond what can be stored. This
exponential growth is the result of the curse of dimensionality.

\subsection{PDEs with Constant Coefficients}
\begin{definition}[Constant-coefficient PDE]

  A linear second-order elliptic PDE has constant coefficients
  if it can be written in the form

  \[
    -\sum_{k=1}^d c_k \frac{\partial^2 u}{\partial x_k^2}
    = f(x),
  \]

  where the coefficients $c_k \in \mathbb{R}$ are constants (independent of
  $x$).

  In the special case $c_k = 1$ for all $k$, this reduces to the Poisson
  equation

  \[
    -\Delta u = f.
  \]

\end{definition}
\paragraph{Discretized form.}

After discretization on a tensor grid with $n$ points per dimension,
the PDE becomes a linear system

\[
  A U = F,
\]

where

\[
  A
  =
  \sum_{k=1}^d
  I \otimes \cdots \otimes
  L
  \otimes \cdots \otimes I,
\]

and $L \in \mathbb{R}^{n \times n}$ is the one-dimensional discrete Laplacian.

\paragraph{Remark.}

The differential operators are applied using the TTD-based
procedures described in \cref{sec:differential_operators}.

\paragraph{Explanation (informal).}

When coefficients are constant, each spatial direction contributes the same type
of operator.

This leads to a separable Kronecker structure in the discretized matrix $A$.

As a result:

\begin{itemize}
  \item $A$ has low TT-rank,
  \item it can be stored efficiently in TT format,
  \item high-dimensional problems become tractable.
\end{itemize}

The structure of the operator does not depend on position, which is what enables
this strong separability.

\subsection{PDEs with Variable Coefficients}
\begin{definition}[Variable-coefficient PDE]

  A second-order elliptic PDE has variable coefficients if it has the form

  \[
    -\nabla \cdot (a(x)\nabla u(x)) = f(x),
  \]

  where the coefficient function

  \[
    a : \Omega \to \mathbb{R}
  \]

  depends on the spatial variables.
\end{definition}
\paragraph{Discretized form.}

After discretization, we again obtain a linear system

\[
  A(a) U = F,
\]

but the matrix $A(a)$ now depends on the spatially varying coefficient.

\paragraph{Explanation (informal).}

When $a(x)$ varies in space, the operator is no longer a simple Kronecker sum.

However, if the coefficient function admits a low-rank tensor approximation, for
example

\[
  a(x_1,\dots,x_d)
  \approx
  \sum_\alpha
  a_1(x_1)\cdots a_d(x_d),
\]

then the discretized operator $A(a)$ can still be approximated in TT format.

Thus, TT methods extend beyond constant-coefficient problems.

\subsection{Structural Lattice of PDEs}

\begin{definition}[Structured PDE components]

A general linear PDE can be written as

\[
  \mathcal{L}u
  =
  -\nabla \cdot (K(x)\nabla u)
  + b(x)\cdot \nabla u
  + c(x)u
  = f(x),
\]

where

\begin{itemize}
  \item $K(x) \in \mathbb{R}^{d \times d}$ is the diffusion matrix,
  \item $b(x) \in \mathbb{R}^d$ is the advection (transport) field,
  \item $c(x)$ is a reaction term,
  \item $f(x)$ is the inhomogeneity.
\end{itemize}

Each component may be:
\begin{itemize}
  \item constant or variable,
  \item scalar or vector/matrix-valued,
  \item separable or non-separable.
\end{itemize}

\end{definition}
We can organize PDEs by increasing structural complexity:

\[
\text{constant coefficients}
\;\subset\;
\text{separable variable coefficients}
\;\subset\;
\text{fully variable coefficients}.
\]

\paragraph{Explanation (informal).}

A PDE is not just ``one object'' — it is composed of several building blocks:
diffusion, transport, reaction, and forcing.

Each of these components can vary in complexity. This naturally induces a
\emph{hierarchy (lattice)} of PDE classes.

%==================================================
\section{Implicit Methods for Solving PDE}
%==================================================

\subsection{Implicit Method}

\begin{definition}[Implicit time discretization]

  Consider a time-dependent PDE of the form

  \[
    \frac{\partial u}{\partial t} = \mathcal{L} u + f,
    \qquad u(0) = u_0.
  \]

  Let $t_n = n\Delta t$.

  An implicit time-stepping scheme computes $u^{n+1}$
  by solving an equation involving the unknown future value:

  \[
    \Phi(u^{n+1}, u^n) = 0.
  \]

  In general form:

  \[
    u^{n+1}
    =
    u^n
    +
    \Delta t \, \mathcal{L}(u^{n+1})
    +
    \Delta t \, f^{n+1}.
  \]

\end{definition}
\paragraph{Explanation (informal).}

In an implicit method, the new solution $u^{n+1}$ appears on both sides of the
equation. This means we must solve a system of equations at every time step.

Advantages:

\begin{itemize}
  \item Stable for large time steps
  \item Often unconditionally stable
  \item Suitable for stiff PDEs
\end{itemize}

Disadvantages:

\begin{itemize}
  \item Requires solving linear systems
  \item Computationally more expensive per step
\end{itemize}

\subsection{Method A}

%==================================================
\section{Explicit Methods for Solving PDE}
%==================================================
\subsection{Explicit Method}

\begin{definition}[Explicit time discretization]

  An explicit time-stepping scheme computes $u^{n+1}$
  using only already known quantities:

  \[
    u^{n+1}
    =
    u^n
    +
    \Delta t \, \mathcal{L}(u^n)
    +
    \Delta t \, f^n.
  \]

\end{definition}
\paragraph{Explanation (informal).}

In an explicit method,
the new value $u^{n+1}$ is computed directly
from previously known values.

Advantages:

\begin{itemize}
  \item No linear system to solve
  \item Cheap per time step
  \item Easy to implement
\end{itemize}

Disadvantages:

\begin{itemize}
  \item Stability restrictions on $\Delta t$
  \item CFL condition
  \item Often requires very small time steps
\end{itemize}


\subsection{Explicit finite difference method}

\paragraph{Applicable PDEs.}

This method is intended for homogeneous advection--diffusion equations
with constant coefficients that can be written in the form
\[
u_t
+
\nabla\cdot(\mathbf a\,u)
+
\nabla\cdot(\mathbf B\,\nabla u)
=
0,
\]
or equivalently
\[
u_t
=
-
\left(
\nabla\cdot(\mathbf a\,u)
+
\nabla\cdot(\mathbf B\,\nabla u)
\right).
\]

Here $\mathbf a$ is the advection vector and $\mathbf B$ is the diffusion
matrix. In our implementation, this corresponds to the homogeneous
vector-advection matrix-diffusion PDE.

\paragraph{Method.}

Using Forward Euler time integration,
\[
u^{n+1}
=
u^n
-
\Delta t
\left(
\nabla\cdot(\mathbf a\,u^n)
+
\nabla\cdot(\mathbf B\,\nabla u^n)
\right).
\]

\paragraph{Advection term.}

The advection contribution is

\[
\alpha
=
\nabla\cdot(\mathbf a\,u).
\]

The flux

\[
\mathbf F_a
=
\mathbf a\,u
\]

is constructed first, after which the divergence operator
is applied:

\[
\alpha
=
\nabla\cdot \mathbf F_a.
\]

\paragraph{Diffusion term.}

The diffusion contribution is

\[
\beta
=
\nabla\cdot(\mathbf B\,\nabla u).
\]

The gradient is computed first:

\[
\mathbf g
=
\nabla u.
\]

The diffusion flux is then

\[
\mathbf F_d
=
\mathbf B\,\mathbf g,
\]

and the divergence yields

\[
\beta
=
\nabla\cdot\mathbf F_d.
\]

\paragraph{Implementation note.}
The gradient and divergence operators are computed using the
TTD-based differential operators described in
\cref{sec:differential_operators}.

\paragraph{Time update.}

The explicit time derivative is

\[
u_t
=
-(\alpha+\beta).
\]

The Forward Euler update is therefore

\[
u^{n+1}
=
u^n
-
\Delta t
(\alpha+\beta).
\]

\begin{algorithm}[ht]
  \caption{Explicit finite difference time step}
  \label{alg:explicit_fd}
  \begin{algorithmic}[1]

    \Require State $u^n$, advection vector $\mathbf a$,
    diffusion matrix $\mathbf B$, time step $\Delta t$

    \State Compute advection flux
    \[
      \mathbf F_a = \mathbf a\,u^n
    \]

    \State Compute advection term
    \[
      \alpha = \nabla\cdot\mathbf F_a
    \]

    \State Compute gradient
    \[
      \mathbf g = \nabla u^n
    \]

    \State Compute diffusion flux
    \[
      \mathbf F_d = \mathbf B\,\mathbf g
    \]

    \State Compute diffusion term
    \[
      \beta = \nabla\cdot\mathbf F_d
    \]

    \State Compute time derivative
    \[
      u_t = -(\alpha+\beta)
    \]

    \State Update solution
    \[
      u^{n+1}
      =
      u^n + \Delta t\,u_t
    \]

  \end{algorithmic}
\end{algorithm}

\paragraph{Stability remark.}

The Forward Euler method is conditionally stable.
The admissible time step depends on the spatial discretization,
the advection velocity, and the diffusion coefficients. For diffusion-dominated problems, the allowable time step may
decrease quadratically with the spatial resolution.

In practice, the time step must satisfy a suitable CFL-type
stability condition.

%==================================================
\section{Conclusion and Further Reading}
%==================================================
% Add citations and key references here (TT-SVD, Oseledets, AMEn, etc.)

\bibliographystyle{alpha}
\bibliography{references}

\end{document}
